% ==== P2 draft o3 — §2.2 대체 블록 (L106–L219 대체) ====
\subsection{단일 삽입 자리 --- 점유 통계와 그 요동}\label{sec:sm-site}
앞 소절이 세운 대정준 구도(식~\eqref{eq:sm-gc})의 최소 단위는 삽입 자리 \emph{하나}다. 이 소절은 그 한
자리를 풀어 평균 점유율 $\theta$ 와 그 요동을 얻는다 --- 내부 자유도 없는 국소화 흡착의 원형을 먼저 세우고,
점유된 이온의 진동 자유도를 얹어 유효 자리 자유에너지로 넓힌 뒤, 그 결과를 다음 소절의 다자리
lattice gas 와 본론 평형 peak(\S\ref{sec:eqpeak})의 입력으로 넘긴다. 자리는 전해질 너머의 Li 저장조와
입자를 교환하므로 식~\eqref{eq:sm-gc} 의 대정준 설정 그대로다.

\textbf{(a) 출발 --- 두 가정과 두-상태 미시상태.} 유도에 앞서 두 가정을 명시한다.
\begin{itemize}[leftmargin=2.2em,itemsep=0.1em]
\item \textbf{가정 1 (hard-core 배타).} 한 자리에는 Li 이온이 $0$ 개 또는 $1$ 개만 들어간다. 근거는 양자
통계가 아니라 강한 단거리 반발이다 --- 자리의 부피와 이온 사이 정전$\cdot$입체 반발이 이중 점유의
에너지를 사실상 무한대로 올린다. 이 가정이 점유수 합을 $n\in\{0,1\}$ 로 자른다.
\item \textbf{가정 2 (자리 독립).} 이 소절에서는 자리의 에너지가 이웃 자리의 점유와 무관하다고 둔다.
이 가정은 \S\ref{sec:sm-mf} 에서 상호작용 파라미터 $\Omega$ 로 완화된다.
\end{itemize}
가장 단순한 기술에서 자리는 두 미시상태만 갖는다 --- 빈 자리($n=0$, 에너지 $0$ 기준)와 점유
자리($n=1$, 에너지 $\varepsilon_0$)다. 점유 이온의 내부 진동은 잠시 접어 두고, 이 두-상태 셈법으로
표준 결과를 먼저 세운다.

\textbf{(b) 연산 --- 대정준 합.} 식~\eqref{eq:sm-gc} 의 대정준 합을 이 두 상태에 쓰면, 각 상태에 Gibbs
인자 $e^{-\beta(\varepsilon_0-\mu)n}$ 을 곱해 더한 두 항이다:
\begin{equation}
\Xi_1^0\;=\;\sum_{n=0,1}e^{-\beta(\varepsilon_0-\mu)\,n}\;=\;1+e^{-\beta(\varepsilon_0-\mu)}.
\label{eq:sm-bareXi}
\end{equation}

\textbf{(c) 중간식 --- 평균 점유.} 자리의 평균 점유는 두 상태의 확률 가중 평균이다:
\begin{equation}
\langle n\rangle_0=\frac{0\cdot1+1\cdot e^{-\beta(\varepsilon_0-\mu)}}{\Xi_1^0}
=\frac{e^{-\beta(\varepsilon_0-\mu)}}{1+e^{-\beta(\varepsilon_0-\mu)}},
\label{eq:sm-baremid}
\end{equation}
식~\eqref{eq:sm-gc} 의 셋째 식 $k_BT\,\partial\ln\Xi_1^0/\partial\mu$ 로 교차검증해도 같은 값이다.

\textbf{(d) 박스 --- 원형 점유율.} 분자$\cdot$분모를 $e^{-\beta(\varepsilon_0-\mu)}$ 로 나누면 점유율이
단일 준위 에너지 $\varepsilon_0-\mu$ 의 logistic 으로 닫힌다:
\begin{equation}
\boxed{\;\Xi_1^0\;=\;1+e^{-\beta(\varepsilon_0-\mu)}\;,\qquad
\langle n\rangle_0\;=\;\frac{1}{1+e^{+\beta(\varepsilon_0-\mu)}}\;.}
\label{eq:sm-bare}
\end{equation}
이 유도는 국소화 흡착(Langmuir 등온선)의 표준 대정준 유도이며 \cite{hill1960,fowler1939}, 삽입 전극에
대한 적용의 원형은 McKinnon--Haering \cite{mckinnon1983} 이다. 식~\eqref{eq:sm-bare} 의
$\langle n\rangle_0$ 은 구조상 Fermi--Dirac 함수와 \emph{동형}이다 --- 단일 준위 에너지가
$\varepsilon_0-\mu$ 인 2-상태(채움/빔) 점유가 전자 준위의 점유/공위와 같은 대수 구조를 낳는다.
공유되는 것이 그 대수 구조뿐이고 배타의 물리적 근거는 서로 다르다는 점은 식~\eqref{eq:fermifn} 아래
검산에서 짚는다.

\begin{bgbox}
\textbf{페르미온$\cdot$보손과 양자 통계 ---} 같은 종류의 입자는 원리적으로 서로 구별되지 않는다. 두
입자를 맞바꾸는 연산 아래 다입자 파동함수가 어떻게 변하는가로 자연의 입자는 두 부류로 갈린다 --- 파동함수가
부호를 유지하면(대칭) \emph{보손}(boson), 부호를 뒤집으면(반대칭) \emph{페르미온}(fermion)이다.
스핀--통계 정리로 정수 스핀 입자가 보손, 반정수 스핀 입자가 페르미온이며, 광자와 포논이 앞의 예,
전자가 뒤의 예다. 보손은 한 양자 상태에 임의로 많이 쌓일 수 있어 평균 점유수가 Bose--Einstein 분포
$\bar n=1/(e^{\beta(\varepsilon-\mu)}-1)$ 를 따른다. 페르미온은 Pauli 배타 원리로 한 양자 상태에 최대
하나만 들어가므로 준위당 점유가 $0$ 또는 $1$ 로 잘리고, 그 평균이 Fermi--Dirac 분포
$f(\varepsilon)=1/(e^{\beta(\varepsilon-\mu)}+1)$ 다. 리튬 삽입 자리의 점유도 $0$ 또는 $1$ 로 잘리지만,
그 배타의 근거는 양자 통계가 아니라 기하다 --- 자리의 부피와 이온 사이 정전$\cdot$입체 반발이 이중 점유를
금한다. ``한 대상에 $0$ 또는 $1$''이라는 셈법의 구조가 페르미온 준위와 공통이어서, 통계의 기원이 전혀
다른데도 같은 대수형(logistic, 곧 Fermi--Dirac 꼴)이 나온다. 진짜 양자 통계는 Chapter 2 에서 실제로
쓰인다 --- 격자 진동의 양자인 포논은 보손이라 Bose--Einstein 분포로, 전도 전자는 페르미온이라
Fermi--Dirac 분포로 각각의 엔트로피를 세운다 \cite{mcquarrie1976}.
\end{bgbox}

\textbf{(a) 출발 --- 점유 상태의 내부 자유도.} 점유수 $n$ 만으로는 미시상태가 다 지정되지 않는다.
$n=1$ 인 자리에서 이온은 한 점에 얼어붙어 있는 것이 아니라 자리 우물 안에서 진동하는 연속 자유도
(위치$\cdot$운동량)를 가진다. 미시상태의 온전한 목록은 ``빈 자리($n=0$, 에너지 $0$ 기준)'' 하나와,
``점유 자리($n=1$, 우물 바닥 에너지 $\varepsilon_0$)에 내부 배치 $\Gamma$ 가 붙은 연속족''이다 --- 두
미시상태 가운데 점유 상태를 그 내부 배치에 대해 펼친 것이다.

\textbf{(b) 연산 --- 내부 배치 적분과 $q(T)$.} 식~\eqref{eq:sm-gc} 의 대정준 합을 이 넓힌 목록
그대로 쓰면, 점유수에 대한 합(가정 1 로 두 항)과 각 점유 상태의 내부 배치에 대한 적분으로 갈라진다:
\begin{equation}
\Xi_1\;=\;\sum_{n=0,1}e^{\beta\mu n}\,z_n
\;=\;1+q(T)\,e^{-\beta(\varepsilon_0-\mu)},
\qquad
q(T)\equiv\int e^{-\beta H_\mathrm{int}(\Gamma)}\,\frac{\dd\Gamma}{h^3}
\label{eq:partfn}
\end{equation}
--- $z_n$ 은 점유수 $n$ 인 상태들의 canonical 분배함수로, $z_0=1$(빈 자리는 상태 하나),
$z_1=e^{-\beta\varepsilon_0}q(T)$(우물 바닥 인자에 내부 자유도의 적분 $q(T)$ 가 곱해진다)이다.
$q(T)$ 는 점유된 이온의 진동 등 국소 자유도의 단일 자리 분배함수다\footnote{이 $q(T)$(단일 자리 내부 분배함수)는 N6--N9 의 정규화 용량 좌표 $q{=}Q/Q_\cell$ 와 다른 양이다 --- 전자는 Part 0 의 국소 열역학량, 후자는 곡선 좌표.}. 일반적인 정의는 상태
합 $q(T)=\mathrm{Tr}\,e^{-\beta H_\mathrm{int}}$ 이고 식~\eqref{eq:partfn} 의 위상공간 적분은 그 고전
극한이다 --- 자리 우물을 3차원 조화 퍼텐셜(진동수 $\omega_0$)로 근사하면 고전 극한에서
$q=(k_BT/\hbar\omega_0)^3$ 이고, 등방 가정(세 방향이 같은 $\omega_0$)을 풀어 데카르트 축마다 다른
진동수를 허용하는 일반형으로 넓히면 양자 상태 합으로는
$q=\prod_{i=1}^{3}[2\sinh(\beta\hbar\omega_i/2)]^{-1}$(축 지표 $i=1,2,3$ 마다 진동수 $\omega_i$)로
닫힌다 \cite{hill1960}. 기호에 관하여 --- $\Xi_1$ 은 자리 1개의 \emph{대정준} 분배함수로,
\S\ref{sec:sm-ensemble} 의 canonical $Z$ 와는 앙상블이 다르므로 기호부터 구분해 적는다(첨자 $1$ = 자리
1개; 내부 자유도를 끈 원형은 $\Xi_1^0$, 식~\eqref{eq:sm-bare}). Chapter 2 의 단일 자리
분배함수(그 문건 식 (eq:Z1))도 같은 기호 $\Xi_1$ 로 통일한다.

\textbf{(c) 중간식 --- 평균 점유와 계수 $q(T)$ 의 흡수.} 점유 상태들의 확률 합이 곧 평균 점유다:
\begin{equation}
\langle n\rangle=0\cdot P_0+1\cdot P_1
=\frac{q(T)\,e^{-\beta(\varepsilon_0-\mu)}}{1+q(T)\,e^{-\beta(\varepsilon_0-\mu)}},
\label{eq:sm-occmid}
\end{equation}
그리고 식~\eqref{eq:sm-gc} 의 셋째 식으로 교차검증하면 $k_BT\,\partial\ln\Xi_1/\partial\mu$ 가 같은
결과를 준다. 여기서 적분이 남긴 계수 $q(T)$ 는 \emph{분자와 분모의 점유-상태 항 양쪽에 같은 인수로}
붙어 있다. 따라서 결과는 $q$ 와 $\varepsilon_0$ 을 따로 아는 데 의존하지 않고 한 조합에만 의존하며,
그 조합은 지수 안으로 흡수된다:
\begin{equation}
q(T)\,e^{-\beta\varepsilon_0}\;=\;e^{-\beta\tilde\varepsilon(T)},
\qquad
\tilde\varepsilon(T)\;\equiv\;\varepsilon_0-k_BT\ln q(T)
\label{eq:sm-epstilde}
\end{equation}
--- $\tilde\varepsilon$ 는 점유 자리의 \emph{자유에너지}(우물 바닥 에너지 $+$ 내부 자유도의 자유에너지
$-k_BT\ln q$)다. 자리 점유의 자유에너지 비용을 $\Delta\mu\equiv\tilde\varepsilon-\mu_\mathrm{Li}$ 로
줄여 쓰면 식~\eqref{eq:sm-occmid} 는 $e^{-\beta\Delta\mu}/(1+e^{-\beta\Delta\mu})$ 가 되어, 내부
자유도가 있어도 함수형은 두-상태 점유의 logistic 그대로 보존된다.

\textbf{(d) 박스 --- 유효 점유율 $\theta$.} 분모$\cdot$분자를 $e^{-\beta\Delta\mu}$ 로 나누면
\begin{equation}
\boxed{\;\theta\;\equiv\;\langle n\rangle\;=\;\frac{1}{1+e^{+\beta\,\Delta\mu}}\;,
\qquad \Delta\mu\equiv\tilde\varepsilon(T)-\mu_\mathrm{Li}\;.}
\label{eq:fermifn}
\end{equation}
\begin{verifybox}
극한 검산 --- $\beta\to\infty$($T\to0$)에서 $\Delta\mu\gtrless0$ 에 따라 $\theta\to0/1$ 의 계단
($T\to0$ 에서 $\tilde\varepsilon$ 는 점유 상태의 바닥 에너지로 수렴한다 --- 고전 $q$ 면 $k_BT\ln q\to0$
이라 $\varepsilon_0$, 양자 조화 $q$ 면 영점 에너지를 더한 $\varepsilon_0+\tfrac32\hbar\omega_0$ --- 어느
쪽이든 상수라 계단 구조는 동일),
$\beta\to0$($T\to\infty$)에서는 $\beta\Delta\mu\to-\ln q$ 이므로 내부 자유도가 없는 기준($q\equiv1$)에서
$\theta\to\tfrac12$(두 상태 등확률)이고, $q(T)$ 가 온도와 함께 증가하는 실제 조화 우물에서는 점유 상태의
위상공간 증가가 이겨 $\theta\to1$ 로 향한다; $\mu_\mathrm{Li}\to\pm\infty$ 에서 $\theta\to1/0$.
\emph{원형 회수} --- 내부 자유도가 없는 극한 $q\equiv1$ 에서 $\tilde\varepsilon\to\varepsilon_0$ 이라
$\Delta\mu\to\varepsilon_0-\mu_\mathrm{Li}$ 이므로, $\mu_\mathrm{Li}=\mu$ 로 읽으면 식~\eqref{eq:fermifn}
은 원형 박스~\eqref{eq:sm-bare} 의 $\langle n\rangle_0$ 를 그대로 회수한다. ★함수형 가드 ---
식~\eqref{eq:fermifn} 은 Fermi--Dirac 함수와 \emph{동형}이지만
물리량은 다르다: 공유하는 것은 ``한 자리에 입자 0 또는 1''이라는 배타 점유의 대수 구조뿐이고(가정 1),
여기의 입자는 전자가 아니라 Li 이온이며 배타성의 근거도 양자 통계가 아니라 자리 하나에 이온 하나라는
기하다.
\end{verifybox}

한 걸음 더 --- $n\in\{0,1\}$ 라 $n^2=n$ 이므로 점유의 요동(분산)이 닫힌 꼴로 나온다:
\begin{equation}
\mathrm{var}(n)=\langle n^2\rangle-\langle n\rangle^2=\theta(1-\theta),
\qquad
\frac{\partial\theta}{\partial(\beta\mu)}=\theta(1-\theta)
\label{eq:sm-flucres}
\end{equation}
(둘째 식은 식~\eqref{eq:fermifn} 의 $\beta\mu$ 미분). \emph{점유의 $\mu$-감도 $=$ 점유 요동} --- 요동--응답
(fluctuation--response) 관계의 최소 사례이고, 본론 평형 peak 의 종 $\xi_\eq(1-\xi_\eq)$(\S\ref{sec:eqpeak})의
통계적 정체가 바로 이 분산이다.

유효 자리 자유에너지 $\tilde\varepsilon(T)$ 가 하류에서 가는 곳은 두 갈래다. 첫째, 온도를 고정하고 보면
$\tilde\varepsilon$ 는 상수 하나이고, \S\ref{sec:sm-electro} 의 전기화학 결선에서 전이 중심 $U_j$ 로
흡수된다 --- 그 소절과 본론이 쓰는 몰당 기준 에너지 $\mu^0$ 는 정확히 $\tilde\varepsilon$ 의 몰 환산
$N_A\tilde\varepsilon$ 이다. 둘째, 온도를 미분하면 내부 자유도의 자유에너지 몫
$f_\mathrm{int}=-k_BT\ln q$ 가 자리당 엔트로피
\begin{equation}
s_\mathrm{int}\;=\;-\frac{\partial f_\mathrm{int}}{\partial T}
\;=\;k_B\,\frac{\partial\,(T\ln q)}{\partial T}
\label{eq:sm-sint}
\end{equation}
를 주는데, 이것이 Chapter 2 의 vibrational 엔트로피 사슬의 단일 자리 원형이다 --- 여기 자리 하나의
세 축 진동수 $\omega_i$($i=1,2,3$)를 격자 전체의 모든 진동 모드를 헤아리는 지표 $k$ 로 넓혀, 조화 모드의
$q_k=[1-e^{-\beta\hbar\omega_k}]^{-1}$ 를 식~\eqref{eq:sm-sint} 에 넣으면 그 문건의 모드당 엔트로피
$s_k=k_B[(1+\bar n_k)\ln(1+\bar n_k)-\bar n_k\ln\bar n_k]$(Bose--Einstein 들뜸수 $\bar n_k$)가 그대로
나온다. 기준점에 관한 사소한 주의로, 위 (b) 의 $[2\sinh(\beta\hbar\omega/2)]^{-1}$ 은 우물 바닥 기준
(영점 에너지를 $q$ 에 포함)이고 여기의 $[1-e^{-\beta\hbar\omega}]^{-1}$ 은 영점 에너지를
$\varepsilon_0$ 쪽에 포함시킨 기준이다 --- 두 규약의 $q$ 는 인자 $e^{-\beta\hbar\omega/2}$(영점 에너지
$\hbar\omega/2$ 의 Boltzmann 인자)만큼 다른데, 이 인자가 $T\ln q$ 에 남기는 기여는 온도 무관 상수
$-\hbar\omega/(2k_B)$ 이므로 온도 미분, 곧 엔트로피는 동일하다. 따라서 전이 중심의 온도 기울기
$\dd U_j/\dd T=\Delta S_{\rxn,j}/F$(\S\ref{sec:center})에서 진동
엔트로피가 차지하는 몫의 미시적 기원이 바로 $\tilde\varepsilon(T)$ 의 온도 의존, 곧 $q(T)$ 다.
% ==== 블록 끝 ====
% NOTE:
% Read Coverage: ch1_sec02a_part0.tex 전문 L1–269(교체 대상 L106–219 포함); V1020_STYLE_RUBRIC.md 전문 L1–51;
%   ch1_sec02b_part0.tex L1–60(§2.4 가 받는 μ(θ)·eq:sm-Smix·eq:mu 접합); ch2_sec01_partition.tex L16–50(eq:Z1·eq:occ
%   원형-선행 서술·FD 동형·Ξ₁ 통일·s_int 원형); ch1_preamble.tex L24–43(\newtheorem*{bgbox}{배경} 확인); ch1_bib.tex
%   L1–30(인용키 실재 검증); V1020_REFERENCE_LEDGER.md(ashcroftmermin1976 = Ch2 스코프 확인).
% 보존 체크:
%   P-1 가정 1/2 = 원형 (a) itemize(hard-core 배타·자리 독립, \S sec:sm-mf 완화).
%   P-2 미시상태 온전(점유 상태 연속 내부 배치 Γ) = 확장 (a) "온전한 목록 … 내부 배치 Γ 붙은 연속족".
%   P-3 라벨 전존·수식 불변 = eq:partfn(확장 b)·eq:sm-occmid(확장 c)·eq:sm-epstilde(확장 c)·eq:fermifn(확장 d)·
%       eq:sm-flucres(요동 단락)·eq:sm-sint(둘째 갈래) — 여섯 최종 수식 원문 그대로.
%   P-4 q(T)↔용량 좌표 q 각주 = 확장 (b) eq:partfn 직후 \footnote 원문 보존.
%   P-5 조화 우물 q 고전 (k_BT/ħω₀)³·양자 ∏[2sinh(βħω_i/2)]^{-1} = 확장 (b) 원문 보존.
%   P-6 ★FD 동형 가드 = verifybox 말미(공유는 대수 구조뿐·입자는 Li 이온·근거는 기하) — 자족 서술, bgbox 없이도 성립.
%   P-7 요동-응답 var(n)=θ(1-θ)=∂θ/∂(βμ) = eq:sm-flucres 및 직후 문장.
%   P-8 ε̃ 두 갈래(U_j 흡수·μ⁰=N_Aε̃ / s_int=k_B∂(T ln q)/∂T·Ch2 BE s_k·영점 에너지 규약 단락·dU_j/dT) = 후미 2문단 원문 보존.
%   P-9 grand canonical 문맥(eq:sm-gc 셋째 식 교차검증) = 원형 (c) k_BT∂lnΞ₁⁰/∂μ + 확장 (c) k_BT∂lnΞ₁/∂μ 이중.
%   P-10 Ch2 기호 통일(Ξ₁) = 확장 (b) 말미 "Chapter 2 … 같은 기호 Ξ₁ 로 통일"(+Ξ₁⁰ 원형 구분 명시).
% 신규 라벨: eq:sm-bareXi(원형 Ξ₁⁰ 합)·eq:sm-baremid(원형 ⟨n⟩₀ 중간식)·eq:sm-bare(원형 박스 Ξ₁⁰·⟨n⟩₀).
% 자체검수:
%   (극한/부호) 원형 박스는 β→∞ 계단·FD 동형만 언급, 확장 verifybox 가 β→∞/β→0/μ→±∞ 4검산+q≡1 원형 회수(eq:sm-bare)
%     명시 — D7 회수 검산 충족. eq:fermifn 부호 e^{+βΔμ} 원문과 일치, q≡1·μ_Li=μ 대입 시 ⟨n⟩₀ 부호 정합 확인.
%   (rubric) D7 2단(원형 박스 eq:sm-bare→bgbox→확장 박스 eq:fermifn→회수 검산) 분리 준수; D1/D2 자기 diff·버전 서술 0
%     (roadmap·"원형/원형 회수"는 물리 원형어로 원문 관용 계승, 개정 서술 아님); C3 페르미온/보손·FD/BE 하이픈 준수;
%     bgbox 자족(본문 지시대명사 무·μ,ε 일반기호)·mcquarrie1976 만 인용(ashcroftmermin1976 = ch1_bib 미등재 → 미인용, 브리프 "필요시" 재량).
